% =====================================================================
\section{Objetivos}
\label{sec:objetivos}

\subsection{Objetivo geral}

Investigar e estender metodologias de verificação formal baseadas em SMT, por meio do
verificador ESBMC, para componentes de inteligência artificial generativa e de redes
neurais quantizadas, de modo a obter garantias de segurança, limites operacionais estritos
e corretude lógica tanto em abstrações de alto nível quanto em \textit{kernels} de baixo
nível.

\subsection{Objetivos específicos e grau de cumprimento}

O plano de trabalho previa cinco frentes, organizadas por complexidade arquitetural
crescente. A Tabela~\ref{tab:objetivos} apresenta cada uma delas ao lado do que foi
efetivamente entregue, indicando explicitamente onde o resultado ficou aquém do previsto e
por quê. O detalhamento está na Seção~\ref{sec:resultados}.

\begin{table}[H]
\centering
\caption{Objetivos específicos do plano de trabalho e grau de cumprimento.}
\label{tab:objetivos}
\small
\begin{tabular}{@{}p{4.3cm}p{7.1cm}p{2.5cm}@{}}
\toprule
\textbf{Objetivo específico} & \textbf{Resultado obtido} & \textbf{Situação} \\
\midrule
Verificação de redes neurais quantizadas (simulação QNN) &
Camada Q8.8 implementada e verificada em C; \textit{harness} XOR provado em 0,08\,s;
o \textit{frontend} Python do ESBMC não está disponível no binário utilizado &
Cumprido com ressalva \\
\addlinespace
Segurança de \textit{kernels} e \textit{stubs} de inferência &
Contraexemplo de segurança de memória obtido em 0,43\,s; GEMM com \textit{tiling}
provado até $N{=}3$; barreira de escalabilidade medida em $N{=}4$ &
Cumprido \\
\addlinespace
Verificação de propriedades de sistema em malha fechada &
Planta \textit{cart-pole} e ator DDPG quantizado verificados acoplados; duas
propriedades vácuas identificadas e reescritas; envelope de segurança caracterizado &
Cumprido parcialmente \\
\addlinespace
Arcabouço programático em Python e ciclo neuro-simbólico &
API \texttt{core\_verify} com status tri-estado e suíte de regressão; ciclo
gerar--verificar--corrigir demonstrado com \textit{mock} determinístico &
Cumprido com ressalva \\
\addlinespace
Garantias de segurança em aprendizado por reforço &
Política de RL verificada; \textit{solver} produziu contraexemplo: a ação
ultrapassa o limite do atuador na caixa de estados analisada &
Cumprido (resultado negativo) \\
\bottomrule
\end{tabular}
\end{table}

Três observações são necessárias para que a tabela seja lida corretamente. Primeiro,
``cumprido com ressalva'' na primeira linha significa que o objetivo foi atingido em C, e
não em Python: o binário do ESBMC versionado no repositório (6.8.0) não possui o
\textit{frontend} Python compilado e rejeita arquivos \texttt{.py}; concluir essa frente
exige compilar a versão 8.0.0 com \texttt{ENABLE\_PYTHON\_FRONTEND=ON}. Segundo, na quarta
linha, o ciclo neuro-simbólico foi demonstrado com um agente simulado determinístico, e não
com um modelo de linguagem real --- o experimento comprova o \emph{protocolo}
(gerar, verificar, extrair diagnóstico, corrigir), não a autonomia de um agente. Terceiro, o
resultado negativo da última linha é reportado como resultado, e não como falha de
execução: o \textit{solver} exibiu um traço concreto em que a ação excede o intervalo
$[-1,1]$, de modo que o \textit{safety shield} projetado ainda \textbf{não} está validado
como bloqueio eficaz.
